% !TeX spellcheck = en_us
\documentclass[a4paper,12pt,fleqn]{article}

\usepackage[left=1.50cm, right=1.50cm, top=3cm, bottom=2cm, 
headheight= 2.5cm, headsep= 4pt, footskip=25pt]{geometry}
\usepackage{fancyhdr}
\usepackage{lmodern}
\usepackage[table]{xcolor}
\usepackage{amssymb,amstext,amsmath}
\usepackage{eso-pic}
\usepackage{booktabs}
\usepackage{bm}
\usepackage{setspace}
\usepackage{graphicx}
\usepackage[colorlinks]{hyperref}
\setlength{\parindent}{0mm}
\setlength{\parskip}{2mm}
\usepackage{enumitem}
\usepackage{titlesec}
\usepackage{gensymb}
\usepackage{pgfgantt}
\usepackage{lscape}
\usepackage[absolute]{textpos}
\usepackage{comment}
\usepackage{cite}
\usepackage{tikz}
\usepackage{pgfplots}
\usetikzlibrary{shapes.geometric}
\usepackage{array}
\newcolumntype{P}[1]{>{\centering\arraybackslash}p{#1}}
\usepackage{longtable}
\usepackage[font={small}]{caption}
\usepackage[most]{tcolorbox}
\usepackage{tikz}
\usepackage{stmaryrd}
\usepackage{wrapfig}
%\usepackage[none]{hyphenat}
\usepackage{ctable}
\usepackage{ragged2e}
\usepackage[official]{eurosym}
\newcommand*\circled[1]{\tikz[baseline=(char.base)]{
		\node[shape=circle,draw,inner sep=1pt] (char) {#1};}}
\usepackage[symbol]{footmisc}
\renewcommand{\thefootnote}{\fnsymbol{footnote}}
\renewcommand{\thempfootnote}{\fnsymbol{mpfootnote}}
\renewcommand*\footnoterule{}

\usepackage{xcolor}
\definecolor{docColor}{cmyk}{1.00,1.00,0, 0.40}  % 0.90\% of black
\color{docColor}

\newtcolorbox{myhlbox}[2][]{colback=docColor!05, colframe=docColor!20, width=1.0\linewidth,
	boxrule=0pt, left=1em, right=1em, top=1em, bottom=1em,
	colbacktitle=docColor!20, coltitle=docColor,enhanced, attach boxed title to top left={xshift=2mm, yshift=-1mm},
	title=#2, #1}

\hypersetup{
	colorlinks=true,
	citecolor=blue,
	filecolor=magenta,
	linkcolor=blue,
	urlcolor=blue,
	pdftitle={SeminarAnnouncements},
	%bookmarks=true,,
}

\graphicspath{{./figs/}}

\definecolor{dgreen}{rgb}{0.0, 0.5, 0.0}

\allowdisplaybreaks

\newcommand*\mystrut[1]{\vrule width0pt height0pt depth#1\relax}
\renewcommand{\thesection}{\arabic{section}.}
\renewcommand{\thesubsection}{\arabic{section}.\arabic{subsection}}
\renewcommand{\thesubsubsection}{\arabic{section}.\arabic{subsection}.\arabic{subsubsection}}
\titleformat{\section}{\sffamily\large\bfseries}{\thesection}{1.0em}{}
\titleformat{\subsection}{\sffamily\large\bfseries}{\thesubsection}{1.0em}{}
\titleformat{\subsubsection}{\sffamily\normalsize\bfseries}{\thesubsubsection}{1.0em}{}
\titlespacing\section{0pt}{0pt plus 2pt minus 2pt}{0pt plus 2pt minus 2pt}
\renewcommand{\baselinestretch}{1.00}

\newcommand\CircledImage[1]{%
	\tikz\path[fill overzoom image={#1}, line width=0mm]circle[radius=0.5\linewidth];%
}

\input{format/defs}

\fancypagestyle{titlepage}{
	\fancyhf{}
	\insertLXheader
}

\fancypagestyle{plain}{%
	\fancyhf{}
	\renewcommand{\headrulewidth}{0pt}
	\renewcommand{\footrulewidth}{0pt}
	\insertnormalpage
}

\pagestyle{plain}

\newcommand\TitleBackground{%
	\AddToShipoutPictureBG*{
		\begin{tikzpicture}[remember picture,overlay]
			\node[opacity=0.20, anchor=center] 
			at (current page.center)
			{\includegraphics[height=\paperheight,keepaspectratio]{format/lms-bg.jpg}};
		\end{tikzpicture}
	}
}

\begin{document}
	\sffamily
	\thispagestyle{titlepage}
	\TitleBackground
	\vspace*{-3em}
	\begin{center}
		\huge \textbf{LMS Seminar}
	\end{center}
	\begin{center}
		\large 
		\textbf{From cyclic deformation to fatigue crack nucleation at basal twist grain boundaries in the Ti-6Al-4V titanium alloy} 
	\end{center}
	\begin{center}
		\Large
		Samuel Hemery \\ 
		{\large Institut Pprime, ISAE-ENSMA, Université de Poitiers} \\
		\vspace*{1em}
		{\large \textbf{Date, time, and venue} \\
		January 22, 2026 (2 -- 3 pm), Amphithéâtre Pole Mecanique}
	\end{center}

	\vspace*{-1em}
	
	\begin{myhlbox}{{Abstract}}
		Titanium alloys are widely employed in aerospace applications owing to their outstanding combination of properties. As components often experience cyclic mechanical loadings during service, accurate predictions of fatigue lifetimes are a key milestone for safe and lightweight aircrafts, but require a mechanistic modeling of underlying processes. Recent studies revealed the nucleation of fatigue cracks at basal twist grain boundaries (BTGBs) in $\alpha + \beta$ titanium alloys. The presentation will focus on recent progress in clarifying the deformation and fracture mechanisms operating in these microstructural configurations. In particular, a combination of experimental and simulation techniques, including HR-DIC, EBSD, in-situ mechanical testing, and molecular dynamics simulations, was leveraged to characterize the mechanical behavior at BTGBs in titanium alloys subjected to monotonic and cyclic loadings. The origin of preferential cracking at these microstructure configurations will finally be discussed.
	\end{myhlbox}

\vspace*{1em}

	\begin{myhlbox}{{About the speaker}}
		\begin{wrapfigure}{r}{0.25\textwidth}
			\vspace*{-2em}
			\CircledImage{images/hemery}
		\end{wrapfigure}

		Samuel Hemery is associate professor at Institut Pprime and ISAE-ENSMA (Poitiers, France) since 2014. After engineering studies, he completed a doctoral thesis in materials science (2013) at Ecole Centrale Paris. Early works focused on liquid metal embrittlement of steels in liquid sodium for fast-breeder reactors at CEA. His current research interests are mainly focused on the mechanical behavior of titanium alloys employed in aerospace applications, with an emphasis on fatigue properties. In particular, relationships between microstructure and mechanical properties are studied through the analysis of deformation and fracture mechanisms. A variety of characterization (SEM, EBSD, HR-DIC, in-situ and conventional mechanical testing) and simulation (CP-FFT) techniques are leveraged for this purpose. He also is assistant editor for Metallurgical and Materials Transactions A (Key Reader) since 2017, and qualified to supervise research (HDR) since 2023. 
		
	\end{myhlbox}
\end{document}